Une onde sinusoïdale d'amplitude \qty{\v{A}}{cm}, de longueur d'onde \qty{\v{lam}}{cm} et de fréquence \qty{\v{f}}{Hz} se propage vers les $x$ \textbf{négatifs} sur une corde. À $x = 0$ et $t = 0$, le déplacement de la corde est nul et la vitesse transversale de ce point est positive.

\begin{enumerate}[label=\alph*)]
    \item Écrivez la fonction d'onde $y(x,t)$ sous la forme $y = A\sin(kx \pm \omega t + \phi)$, en donnant les valeurs numériques de $k$, $\omega$ et $\phi$.
    \item Déterminez la vitesse de l'onde.
    \item Quel est le déplacement de la corde à $x = \qty{\v{x1}}{m}$ et $t = \qty{\v{t1}}{s}$ ?
\end{enumerate}
